\subsection{Remaining higher-point and character determinations}

\begin{statement}{BO.1}{Open}{The corrected six- and seven-point problem}
The complete torsion, centralizer-torsor, special-fibre and
character classification requested in task B.1 has not been
determined. The universal unique-frame premise is false at
both six and seven points, by the explicit invariant essential
twists of FP.12.
\end{statement}
A corrected theory must retain a frame as a torsor under its
actual pure centralizer, which in these examples is infinite.
Choosing one element of that torsor does not define a canonical
character. Neither the three rigid five-point classes nor the
arithmetic intersections AF.6 give the missing higher-point
classification. The first undecided extension in this return
is already a coordinate independent of the pure-centralizer
choice for those two involution examples. No list of new
six- or seven-point characters is asserted.

\begin{statement}{BO.2}{Open}{All relative quadratic characters of the gentle group}
For $K_G$ of P-L1, we have not determined
$\operatorname{Hom}_{\rm cont}(K_G,\Ftwo)$, or decided
whether it is generated by $\Delta$.
\end{statement}
B.1 completely classifies the characters that factor through
the proved finite frame pair: they are zero and $\Delta$.
That finite factorization has not been proved for every
relative character. LI.4 constructs the actual relative
characters $\Delta\circ D_P^n$ and proves that $\Delta$
is outside their finite span, but does not prove a positive
iterate nonzero. Therefore it does not yet furnish a second
independent quadratic character. The cubic character on
$\HH$ in FP.33 is a different object: conjugation by complex
conjugation negates its values, so it is not being counted
as a quadratic character or as an ordinary full-group
relative abelian cubic character.

\begin{statement}{BO.3}{Open}{The remaining finite-coefficient and relative-stage homology comparison}
The Alexander and characteristic-zero level-twelve questions
are now resolved in the stated scopes. LI.8 excludes a
$\Delta$-line from the canonical Alexander module, and from
homology of stable finite metabelian covers, in characteristic
different from two. RH.2--RH.3 and LH.1 construct a canonical
finite homology action detecting the full character, with a
sign line in its ordered-basis permutation module.
LH.2 excludes every $\Delta$ constituent from the raw
level-twelve homology in characteristic zero for the full
gentle group. We have not determined all finite-coefficient
constituents of that raw module or completed the precise
identification with every listed stage of LM8 B.31--B.35.
\end{statement}
The distinction between detecting an action character and
containing it as a raw homology constituent is essential.
LH.2 proves the latter nonexistence from an actual
good-reduction Frobenius with no eigenvalue one, using the
genuine arithmetic cover attached to the normalized action.
It does not apply after restricting to $\mathcal H$, and its
determinant need not be invertible in finite coefficients.
The first remaining case is consequently an odd finite
coefficient constituent of the raw level-twelve homology,
with the arithmetic subgroup and full gentle action retained.
B.4 separately
refutes the universal assertion that a module and its
$\Delta$-twist must be nonisomorphic, using
$1\oplus\Delta$. The Galois restrictions are isomorphic
for every module by the proved Galois vanishing of $\Delta$.
